\section{从 benchmark 竞争到问题定义竞争}

姚顺宇认为，2026 年初的前沿模型格局已经不再是简单的“谁追不上谁”。Gemini、OpenAI、Anthropic 这几家在公开 benchmark 上的差距明显收窄，SWE-bench、AIME、IMO 等指标仍有参考价值，但很多时候一两个百分点已经接近噪音，而不是稳定信号。模型的真实差异更多体现在用户体验：Claude 在通用工具使用和 agent 上仍然强，Codex 在纯 coding 方向缩小差距，Gemini 在纯 reasoning 和日常使用上有优势，但这些差异不一定能被单一公开分数捕获。

\begin{figure}[H]
\centering
\includegraphics[width=0.88\textwidth]{frame-competition-models.jpg}
\caption{谈到模型实验室竞争时，访谈把“纸面 benchmark”与“真实用户体验”区分开来。\protect\footnotemark}
\end{figure}
\footnotetext{视频画面时间区间：00:16:20--00:16:55。本图用于标记讨论 benchmark 同质化的上下文。}

\begin{knowledgebox}{读图：benchmark 讨论应看什么}
这张图没有曲线或表格，重点是把读者带回访谈现场：这一段正在从公司竞争切到评估方法。真正要比较的是三层信息：公开榜单、内部评估、真实用户体验；画面只证明该论述来自原视频对应区间。
\end{knowledgebox}


\begin{importantbox}{核心判断：差异从“能不能做”转向“要它怎么做”}
早期模型差异常表现为能力差异：谁更会推理、谁更会写代码、谁更会调用工具。当前前沿模型接近后，难点转向定义目标行为：什么才算好 agent，什么才算好 coding assistant，什么样的 long-horizon 行为值得优化。也就是说，竞争对象从单纯模型能力转向任务定义、数据构造、评估体系和产品反馈闭环。
\end{importantbox}

\subsection{术语消化：公开指标为何会失真}

\noindent\begin{tabularx}{\textwidth}{p{0.21\textwidth}p{0.30\textwidth}X}
\toprule
术语 & 解决的问题 & 访谈中的含义 \\
\midrule
SWE-bench & 衡量模型修复真实软件 issue 的能力 & 当前强模型已接近高分区，一两个百分点差异未必代表真实产品体验差异。 \\
AIME / IMO & 数学推理 benchmark & 能显示 reasoning 能力，但不能直接说明模型在工具使用、产品协作、长期任务中的表现。 \\
Noise vs Signal & 区分随机波动与稳定趋势 & 姚顺宇强调，公开榜单接近时，小差距更可能是噪音；要看内部评估和真实使用。 \\
Behavior Definition & 明确模型应当表现出的行为 & 当前核心难点：不是“模型会不会”，而是“应该怎样做才算对”。 \\
\bottomrule
\end{tabularx}

\begin{warningbox}{不要把榜单分数当成完整能力画像}
榜单通常压缩了任务、数据分布、评分规则和失败模式。当前模型实验室的竞争越来越像“高维行为优化”：同样 80\% 的 SWE-bench 分数，可能对应完全不同的交互体验、工具调用策略、失败恢复能力和用户信任感。
\end{warningbox}

\subsection{为什么意愿曾经决定能力差异}

访谈中反复出现一个词：意愿。姚顺宇说，过去模型能力差异很大程度来自公司愿不愿意把资源投向某类能力。Claude 长期重视 tool use 和 coding，OpenAI 一段时间重视 reasoning，这些选择都会影响 infra、数据、评估和训练管线。前沿模型训练不是只靠一个算法想法，而是靠持续投入把行为目标变成可训练的数据和可追踪的指标。

\begin{knowledgebox}{数据不是中性资源}
姚顺宇举了一个很有意思的例子：早期模型写代码好，部分原因可能不是大家预先设计了多么高级的 coding pipeline，而是互联网中的 GitHub 数据天然比普通网页质量更高。也就是说，模型行为有时来自数据分布中的隐性偏置；研究者事后才理解“为什么它学会了这个”。
\end{knowledgebox}

\subsection{本章小结}

模型竞争进入新阶段后，公开 benchmark 仍重要，但不再足以解释全部差异。前沿实验室真正竞争的是：谁能更好地定义任务、构造数据、搭建评估，并把这些转化成稳定的模型行为。

